\section{Methods}
\subsection{Cohort and Prediction Task}
We use the original observed JRA flat-race cohort, with race results and pedigree records collected through Netkeiba \cite{netkeibaRace,netkeibaDb}. Eligible entries are divided by race date into 2010--2021 training, 2022--2023 validation and 2024--2025 test partitions, retaining complete race rosters (Table~\ref{tab:cohort}). The observed collection is not the entire JRA racing universe \cite{jraStats2yo,jraStats3up}. The binary label indicates an official top-three finish; documented non-finishers are negative, pre-start withdrawals are excluded, and all horses sharing a top-three position in a dead heat are positive. Models are fitted separately for turf and dirt. The test races had been inspected in earlier analyses, making this study exploratory rather than a fresh confirmatory test.

\begin{table}[!t]
\caption{Chronological Cohort and Test Surface Denominators}
\label{tab:cohort}
\centering
\begin{tabular}{lrrr}
\hline
Partition & Years & Races & Entries \\
\hline
Training & 2010--2021 & 38,447 & 548,756 \\
Validation & 2022--2023 & 6,401 & 88,106 \\
Test, turf & 2024--2025 & 3,246 & 43,477 \\
Test, dirt & 2024--2025 & 3,219 & 45,660 \\
\hline
Total & 2010--2025 & 51,313 & 725,999 \\
\hline
\end{tabular}
\end{table}


All four arms use the same 43 race-context and recent-history predictors: 12 context variables, an observed-start count, three history-availability indicators and nine attributes for each of up to three previous starts. History uses eligible observed starts strictly before the target date, including observed jump starts, although targets are flat races. Missing history remains missing. Final odds and popularity are not predictors in these four arms. Category maps are fitted on each training surface, distinguishing missing, unseen and categories occurring fewer than five times; classifiers use native categorical handling and numeric missing values.

\subsection{Graph and Feature Pipelines}
The original three-generation graph is the union of pedigrees rooted at training-period starters, with 62,270 roots, 93,474 nodes and 173,412 edges. Parent--child edges are projected to an undirected graph for Node2Vec \cite{grover2016}. The fixed configuration uses 64 dimensions, 10 walks per node, walk length 80, window 5, five epochs, $p=1$ and $q=0.5$. Each seed has its own fixed vectors and 64 cluster centroids. Within a seed, both graph arms use the same original vector space and centroids. The graph is not reconstructed for each target date or held-out horse. It includes relationships contributed by later training roots, and the focal horse may have influenced graph topology even when its own vector is not directly used. Pedigree records were collected retrospectively, and their availability at each historical cutoff was not verified.

The four feature pipelines are H (context and history); $G_{\mathrm{mixed}}$ (H plus the original graph group, using the focal vector when available and otherwise an ancestor vector); $G_{\mathrm{ancestor}}$ (the same graph group and frozen space, applying the ancestor rule to every entry); and P (H plus nine chronological parent-performance features). The graph group comprises 64 vector coordinates and five descriptors: representation status, selected depth, number of available ancestors, missingness and nearest-centroid category. For the ancestor rule, parents, grandparents and great-grandparents are examined in order. At the first generation containing embedded ancestors, their unique vectors are averaged equally. $G_{\mathrm{ancestor}}$ excludes a focal identifier encountered through a cycle; $G_{\mathrm{mixed}}$ retains the original fallback implementation. If no ancestor is available, the vector is zero with an explicit missing status and cluster. Otherwise the closest frozen centroid is assigned by squared Euclidean distance, with lowest index resolving an exact tie. Table~\ref{tab:focused_coverage} reports coverage. The same $G_{\mathrm{ancestor}}$ rule is used across partitions, but coverage, depth and vector distributions can still differ. The two graph arms compare complete pipelines, not an isolated intervention on fallback.

\begin{table*}[!t]
\caption{Graph Representation Coverage by Partition and Surface}
\label{tab:focused_coverage}
\centering
\begin{tabular}{llrrrrrr}
\hline
& & & $G_{\mathrm{mixed}}$ & \multicolumn{3}{c}{$G_{\mathrm{ancestor}}$: first available depth} & Missing \\
Partition & Surface & Entries & direct & 1 & 2 & 3 & (either) \\
\hline
Training & Turf & 271,549 & 271,549 & 271,549 & 0 & 0 & 0 \\
Training & Dirt & 277,207 & 277,207 & 277,207 & 0 & 0 & 0 \\
Validation & Turf & 42,783 & 21,708 & 42,580 & 202 & 1 & 0 \\
Validation & Dirt & 45,323 & 24,398 & 45,081 & 242 & 0 & 0 \\
Test & Turf & 43,477 & 5,641 & 42,866 & 596 & 15 & 0 \\
Test & Dirt & 45,660 & 5,506 & 44,613 & 1,039 & 8 & 0 \\
\hline
\end{tabular}
\par\smallskip Counts are entries, not distinct horses; counts are identical across graph seeds. $G_{\mathrm{mixed}}$ uses fallback for entries without direct vectors. All test entries have sire, dam and damsire identifiers. Their P zero-prior counts are 13,103/17,941/1,982 on turf and 13,182/20,027/2,371 on dirt, in that order.
\end{table*}


\subsection{Chronological Parent Control}
P adds, for sire, dam and damsire separately, an identifier-availability indicator, a count of eligible prior same-surface starts and a smoothed top-three rate. The count and successes exclude all starts of the focal horse and include only training outcomes strictly earlier than the target date; same-day outcomes are unavailable to each other. For a count $n$ and positive count $s$, the fixed rate is $(s+20\times0.20)/(n+20)$. An absent identifier and a known identifier with no qualifying starts both have count zero and rate 0.20, distinguished by the availability indicator. These are supervised outcome-derived summaries, not genetic-effect estimates; the constant prior is a prespecified regularizer rather than an estimated population rate.

Training summaries expand by date, while the map is frozen at the end of 2021 for all validation and test entries, with focal-horse contributions removed. It never incorporates validation or test outcomes. The source history starts in 2010, so earlier parent performance is left-censored; the frozen map also becomes older relative to later targets. The categorical preprocessing setting that disables external target encoding does not remove P's supervised parent summaries.

\subsection{Historical Market Comparison}
The original analysis on the same cohort is retained separately. A uses context and history; B adds categorical sire/dam identifiers; C adds the original graph group; D adds retrospectively recorded final odds and popularity to A; E adds those market fields to C; and F uses final odds alone. The historical comparisons C--A, C--B, E--D and E--C use independently tuned pipelines. Seven of 60 selected historical tuning curves reached the cap, with best rounds potentially earlier. Final odds and popularity enter together and their individual contributions cannot be separated. The four new arms do not include market fields and do not test $G_{\mathrm{ancestor}}$ or P beyond market information. Final-field, going, weight and market timestamps were not verified at a common pre-start horizon.

\subsection{Fitting and Evaluation}
XGBoost, LightGBM and CatBoost \cite{chen2016,ke2017,prokhorenkova2018} are tuned independently for each arm, family and surface. Four fixed candidates are fitted with seed 0, a maximum of 1,200 rounds and early-stopping patience of 75. Candidate selection uses validation average precision (AP); candidates within $10^{-6}$ of the highest AP are ordered by higher validation Hit@3, fewer best rounds and then candidate identifier. The selected parameters and best-round count are refitted using the training partition alone for seeds 0, 1 and 2; validation is not added to the final fit. Of 96 candidate tuning curves, 21 reached the cap; four of 24 selected curves reached it, while no selected best round or refit round was 1,200. A capped curve can have an earlier best round. Equal budgets do not establish equally adequate tuning. Graph and classifier seeds are paired, so their sources of variation are not separately identified.

The primary score is entry-level AP \cite{saito2015}, which aggregates entries with equal scores at the same threshold. The secondary score, race-mean Hit@3, is the fraction of the three highest-scored entries in each full race that finished in the top three, averaged across races. For this within-race selection only, score ties are ordered by horse number and then horse identifier. We average seed-specific scores and paired differences, not predictions. The primary contrasts are $G_{\mathrm{ancestor}}$--H and $G_{\mathrm{ancestor}}$--P; $G_{\mathrm{ancestor}}$--$G_{\mathrm{mixed}}$ is secondary. All are reported by family and surface for both scores.

For each surface, the same 2,000 whole-race bootstrap draws are used across arms, families and seeds. Each draw resamples complete races with replacement, retains repeated draws, recomputes each seed's score, and averages the paired differences across seeds. The 2.5th and 97.5th percentiles give descriptive paired 95\% intervals conditional on the fitted models. They do not include retraining uncertainty, adjust for multiple comparisons or restore confirmatory inference after prior test examination; repeated horses also create dependence beyond races. An interval containing zero does not establish equivalence.
